% !TeX root = analysis-in-several-complex-variables.tex
In this section, we introduce some basic concepts and results in complex analysis with multiple variables.

\subsection{Holomorphic functions}

    We identify \(\bbC^n \cong \bbR^{2n}\).

    \begin{definition}\label{def:differentiable_and_holomorphic_map}
        A continuous map \(f: \bbR^{2n} \to \bbR^{2m}\) is \emph{differentiable} at \(p \in \bbR^{2n}\) if there exists a linear map \(\upd f_p: \bbR^{2n} \to \bbR^{2m}\) such that 
        \[ f(z) = f(p) + \upd f_p(z - p) + o(|z - p|). \]

        A continuous map \(f: \bbC^n \to \bbC^m\) is \emph{holomorphic} at \(p \in \bbC^n\) if it is differentiable at \(p\) and \(\upd f_p\) is \(\bbC\)-linear, i.e., \(\upd f_p(\im z) = \im \upd f_p(z)\) for all \(z \in \bbC^n\).
    \end{definition}

    By a ``function'', we always mean a complex-valued function, i.e., a map \(f: \bbC^n \to \bbC\).
    Fix a coordinate system \(z = (z_1, \ldots, z_n)\) on \(\bbC^n\) and write \(z_j = x_j + iy_j\) for \(j = 1, \ldots, n\).
    Then a differentiable function \(f = u+iv: \bbC^n \to \bbC\) is holomorphic at \(p\) if and only if the Cauchy-Riemann equations hold:
    \[ \frac{\partial u}{\partial x_i}(p) = \frac{\partial v}{\partial y_i}(p), \quad \frac{\partial u}{\partial y_i}(p) = -\frac{\partial v}{\partial x_i}(p), \quad i = 1, \ldots, n. \]
    
    For convenience, we consider the complexified tangent space \(T\bbR^{2n} \ten_\bbR \bbC\) and introduce the following operators.

    \begin{definition}\label{def:Wirtinger_operators}
        The \emph{Wirtinger operators} are defined as 
        \[ \frac{\partial}{\partial z_j} := \frac{1}{2} \left( \frac{\partial}{\partial x_j} - \im \frac{\partial}{\partial y_j} \right), \quad \frac{\partial}{\partial \bar{z}_j} := \frac{1}{2} \left( \frac{\partial}{\partial x_j} + \im \frac{\partial}{\partial y_j} \right), \quad j = 1, \ldots, n. \]
    \end{definition}

    Then we can rewrite the Cauchy-Riemann equations as
    \[ \frac{\partial f}{\partial \bar{z}_j} = 0, \quad j = 1, \ldots, n. \]

    We summarize some properties of Wirtinger operators in the following proposition.
    \begin{proposition}\label{prop:properties_of_Wirtinger_operators}
        The Wirtinger operators satisfy the following properties:
        \begin{enumerate}
            \item \(\partial_{z_j} z_i = \delta_{ij}\), \(\partial_{z_j} \bar{z}_i = 0\), \(\partial_{z_j} \bar{z}_i = 0\), \(\partial_{z_j} \bar{z}_j = \delta_{ij}\);
            \item \(\overline{\left( \partial_{z_j} f \right)} = \partial_{\bar{z}_j} \bar{f}\);
            \item suppose we have \(\bbC^n \xrightarrow{g} \bbC^m \xrightarrow{f} \bbC^l\) and the coordinate on \(\bbC^m\) is \(w = (w_1, \ldots, w_m)\), then the chain rule holds:
                \begin{align*}
                    \frac{\partial (f \circ g)}{\partial z_j} &= \sum_{k=1}^m \frac{\partial f}{\partial w_k}(g(z)) \frac{\partial g_k}{\partial z_j}(z) + \sum_{k=1}^m \frac{\partial f}{\partial \bar{w}_k}(g(z)) \frac{\partial \bar{g}_k}{\partial z_j}(z), \\
                    \frac{\partial (f \circ g)}{\partial \bar{z}_j} &= \sum_{k=1}^m \frac{\partial f}{\partial w_k}(g(z)) \frac{\partial g_k}{\partial \bar{z}_j}(z) + \sum_{k=1}^m \frac{\partial f}{\partial \bar{w}_k}(g(z)) \frac{\partial \bar{g}_k}{\partial \bar{z}_j}(z).
                \end{align*}
        \end{enumerate}
    \end{proposition}
    \begin{proof}
        By direct computation.
    \end{proof}

    We can also consider the complexified of derivatives
    \[ (\upd f_p)_\bbC: T\bbR^{2n} \otimes_\bbR \bbC \to T\bbR^{2m} \otimes_\bbR \bbC. \]
    If we take \(\{\partial_{z_i}, \partial_{\bar{z}_i}\}_{i=1}^n\) as a basis of \(T\bbR^{2n} \otimes_\bbR \bbC\) and \(\{\partial_{w_j}, \partial_{\bar{w}_j}\}_{j=1}^m\) as a basis of \(T\bbR^{2m} \otimes_\bbR \bbC\),
    then the matrix representation of \((\upd f_p)_\bbC\) is
    \[ (\upd f_p)_\bbC = \begin{bmatrix}
        \partial_z f(p) & \partial_{\bar{z}} f(p) \\
        \overline{\partial_{\bar{z}} f(p)} & \overline{\partial_z f(p)}
    \end{bmatrix}. \]
    In particular, if \(f\) is holomorphic, then we have \(\det (\upd f_p)_\bbC = |\det(\partial_z f)(p)|^2 \geq 0\).

    \begin{definition}\label{def:biholomorphic_map}
        A map \(f: \Omega \to \Omega'\) between two open sets \(\Omega \subset \bbC^n\) and \(\Omega' \subset \bbC^m\) is \emph{biholomorphic} if it is a bijection and both \(f\) and \(f^{-1}\) are holomorphic.
    \end{definition}

    If \(f\) is biholomorphic at \(p\), then \(m = n\) and \(\det \upd f_p > 0\).

    \begin{theorem}[Holomorphic Inverse Function Theorem]\label{thm:holomorphic_inverse_function_theorem}
        Let \(f: \bbC^n \to \bbC^n\) be a holomorphic map. 
        If the Jacobian determinant \(\det \upd f_p\) is nonzero at \(p \in \bbC^n\), then there exist open neighborhoods \(U\) of \(p\) and \(V\) of \(f(p)\) such that \(f: U \to V\) is a biholomorphism.
    \end{theorem}
    \begin{proof}
        By the real inverse function theorem, there exist open neighborhoods \(U\) of \(p\) and \(V\) of \(f(p)\) such that \(g=f^{-1}: V \to U\) is a differentiable map.
        It suffices to show that \(g\) is holomorphic.
        By the chain rule (\cref{prop:properties_of_Wirtinger_operators}), since \(f\) is holomorphic, we have
        \[ 0 = \left(\frac{\partial (f\circ g)_i}{\partial \bar{z}_j}\right)(q) = \left(\frac{\partial f_i}{\partial w_k}\right)(g(q)) \left(\frac{\partial g_k}{\partial \bar{z}_j}\right)(q). \]
        Since \(\det(\partial f/\partial w)(f(q)) \neq 0\), the matrix \((\partial f_i/\partial w_k)(g(q))\) is invertible, which implies that \((\partial g_k/\partial \bar{z}_j)(q) = 0\) for all \(k,j\).
        Thus \(g\) is holomorphic.
    \end{proof}

    \begin{theorem}[Holomorphic Implicit Function Theorem]\label{thm:holomorphic_implicit_function_theorem}
        Let \(f: \bbC^{n+m} \to \bbC^m\) be a holomorphic map. 
        Write the coordinates of \(\bbC^{n+m}\) as \((z,w) = (z_1, \ldots, z_n, w_1, \ldots, w_m) \in \bbC^n \times \bbC^m\).
        If \(\det(\partial f/\partial w) \neq 0\) at \((z_0, w_0) \in \bbC^{n+m}\) with \(f(z_0, w_0) = 0\),
        then there exist open neighborhoods \(U\) of \(z_0\) and \(V\) of \(w_0\), and a unique holomorphic map \(g: U \to V\) such that for any \((z, w) \in U \times V\), \(f(z, w) = 0\) if and only if \(w = g(z)\).
    \end{theorem}
    \begin{proof}
        By real implicit function theorem, there exist differentiable map \(g:U \to V\) satisfying the above condition.
        It suffices to show that \(g\) is holomorphic.
        Let \(G: U \to U \times V\) be defined by \(G(z) = (z, g(z))\).
        Then we have \(f \circ G = 0\).
        By the chain rule, we have
        \[ 0 = \frac{\partial (f\circ G)_i}{\partial \bar{z}_j}(q) = \sum_{k=1}^n \frac{\partial f_i}{\partial w_k}(G(q)) \frac{\partial z_k}{\partial \bar{z}_j}(q) + \sum_{l=1}^m \frac{\partial f_i}{\partial w_l}(G(q)) \frac{\partial g_l}{\partial \bar{z}_j}(q). \]
        Since \(\det(\partial f/\partial w)(G(q)) \neq 0\), the matrix \((\partial f_i/\partial w_k)(G(q))\) is invertible, which implies that \((\partial g_l/\partial \bar{z}_j)(q) = 0\) for all \(l,j\).
        Thus \(g\) is holomorphic.
    \end{proof}

\subsection{Cauchy Integral Formula}

    Recall the Cauchy Integral Formula in one complex variable:

    \begin{theorem}[Cauchy Integral Formula in one complex variable]\label{thm:Cauchy_Integral_Formula_in_one_complex_variable}
        Let \(K \subset \bbC\) be a compact set with piecewise differentiable boundary \(\partial K\), and let \(f\) be differentiable on a neighborhood of \(K\). 
        Then for any \(z\) in the interior of \(K\), we have 
        \[ f(z) = \frac{1}{2\pi \im } \int_{\partial K} \frac{f(\zeta)}{\zeta - z} d\zeta + \frac{1}{2\pi \im} \int_K \frac{\partial f}{\partial \bar{\zeta}}(\zeta) \frac{d\bar{\zeta} \wedge d\zeta}{\zeta - z}. \]
        % \Yang{To be continued...}
    \end{theorem}
    \begin{proof}
        \Yang{By Stokes' theorem. To be continued...}
    \end{proof}

    \begin{theorem}[Cauchy Integral Formula in several complex variables]\label{thm:Cauchy_Integral_Formula_in_several_complex_variables}
        Let \(D \subset \bbC^n\) be a polydisk and \(f\) be holomorphic on a neighborhood of the closure of \(D\). 
        Then for any \(z \in D\),
        \[ f(z) = \frac{1}{(2\pi \im)^n} \int_{\partial D_1 \times \cdots \times \partial D_n} \frac{f(\zeta_1, \ldots, \zeta_n)}{(\zeta_1 - z_1) \cdots (\zeta_n - z_n)} d\zeta_1 \cdots d\zeta_n. \]
        % \Yang{To be continued...}
    \end{theorem}
    \begin{proof}
        \Yang{To be continued...}
    \end{proof}

    \begin{corollary}\label{cor:holomorphic_functions_are_analytic}
        Holomorphic functions are analytic.
        \Yang{To be continued...}
    \end{corollary}

    \begin{proposition}\label{prop:open_mapping_theorem}
        Holomorphic functions are open mappings.
        \Yang{To be continued...}
    \end{proposition}

    \begin{proposition}\label{prop:rigidity_of_holomorphic_functions}
        If a holomorphic function \(f: \Omega \to \bbC\) on a connected open set \(\Omega \subset \bbC^n\) attains its maximum at some point in \(\Omega\), then \(f\) is constant.
        \Yang{To be continued...}
    \end{proposition}

    \begin{lemma}\label{prop:Cauchy_estimates}
        Let \(D \subset \bbC^n\) be a polydisk and \(f\) be holomorphic on a neighborhood of the closure of \(D\). Then for any multi-index \(\alpha = (\alpha_1, \ldots, \alpha_n)\),
        \[ \max_{z \in D} \left| \frac{\partial^{|\alpha|} f}{\partial z_1^{\alpha_1} \cdots \partial z_n^{\alpha_n}}(z) \right| \leq \frac{\alpha!}{r^\alpha} \max_{z \in D} |f(z)|, \]
        where \(r = (r_1, \ldots, r_n)\) is the radius of the polydisk \(D\).
        \Yang{To be continued...}
    \end{lemma}

    \begin{theorem}[Generalized Liouville Theorem]\label{thm:generalized_Liouville_Theorem}
        A holomorphic function \(f: \bbC^n \to \bbC\) on the whole space \(\bbC^n\) that satisfies a polynomial growth condition, i.e., there exist constants \(C > 0\) and \(k \geq 0\) such that 
        \[ |f(z)| \leq C(1 + |z|^k), \quad \forall z \in \bbC^n, \]
        must be a polynomial of degree at most \(k\).
        \Yang{To be continued...}
    \end{theorem}

    \begin{theorem}[Montel's Theorem]\label{thm:Montel's_Theorem}
        A family of holomorphic functions on a domain \(\Omega \subset \bbC^n\) that is uniformly bounded on compact subsets of \(\Omega\) is a normal family, i.e., every sequence in the family has a subsequence that converges uniformly on compact subsets of \(\Omega\) to a holomorphic function or to infinity.
        \Yang{To be continued...}
    \end{theorem}

\subsection{Zero sets of holomorphic functions}

    \begin{theorem}[Hartogs' Extension Theorem]\label{thm:Hartogs'_Extension_Theorem}
        Let \(D \subset \bbC^n\) be a domain with \(n \geq 2\), and let \(K \subset D\) be a compact subset such that \(D \setminus K\) is connected. 
        If \(f: D \setminus K \to \bbC\) is a holomorphic function, then there exists a unique holomorphic function \(\tilde{f}: D \to \bbC\) such that \(\tilde{f}|_{D \setminus K} = f\).
        \Yang{To be continued...}
    \end{theorem}
    \begin{proof}
        \Yang{To be checked}
    \end{proof}

    % \begin{theorem}[Hartogs' Separate Analyticity Theorem]\label{thm:Hartogs'_Separate_Analyticity_Theorem}
    %     Let \(D \subset \bbC^n\) be a domain with \(n \geq 2\), and let \(f: D \to \bbC\) be a function such that for each fixed \(z' = (z_1, \ldots, z_{j-1}, z_{j+1}, \ldots, z_n)\), the function \(f(z', z_j)\) is holomorphic in \(z_j\) for all \(j = 1, \ldots, n\). Then \(f\) is holomorphic on \(D\).
    %     \Yang{To be continued...}
    % \end{theorem}

    \begin{corollary}\label{cor:nonexistence_of_isolated_singularities_in_several_complex_variables}
        In contrast to the one-variable case, isolated singularities do not exist in several complex variables. 
        Specifically, if \(f: D \setminus \{p\} \to \bbC\) is a holomorphic function on a domain \(D \subset \bbC^n\) with \(n \geq 2\) and \(p \in D\), then \(f\) can be extended to a holomorphic function on the entire domain \(D\).
    \end{corollary}
    \begin{proof}
        This is a direct consequence of Hartogs' Extension Theorem by taking \(K = \{p\}\).
    \end{proof}

    \begin{theorem}[Weierstrass Preparation Theorem]\label{thm:Weierstrass_Preparation_Theorem}
        Let \(f: \bbC^{n+1} \to \bbC\) be a holomorphic function in a neighborhood of the origin such that \(f(0) = 0\) and \(f\) is not identically zero. 
        Write the coordinates as \((z,w) = (z_1, \ldots, z_n, w) \in \bbC^{n} \times \bbC\).
        Suppose that \(f(0,w)\) has a zero of order \(k\) at \(w = 0\), i.e.,
        \[ f(0,w) = a_k w^k + a_{k+1} w^{k+1} + \cdots, \quad a_k \neq 0. \]
        Then there exists a neighborhood \(U\) of the origin and unique holomorphic functions \(g: U \to \bbC\) and \(h_j: U' \to \bbC\) for \(j = 1, \ldots, k\), where \(U' \subset \bbC^n\) is the projection of \(U\) onto the first \(n\) coordinates, such that
        \[ f(z,w) = (w^k + h_1(z) w^{k-1} + \cdots + h_k(z)) g(z,w), \]
        with \(g(0) \neq 0\) and \(h_j(0) = 0\) for all \(j\).
        \Yang{To be continued...}
    \end{theorem}
    \begin{proof}
        \Yang{To be continued...}
        \Yang{Use the Cauchy Integral Formula to check the holomorphicity of \(g\) and \(h_j\).}
    \end{proof}

% \subsection{Sheaf of holomorphic functions}

    \begin{definition}\label{def:sheaf_of_holomorphic_functions}
        Let \(\Omega \subset \bbC^n\) be an open set. 
        The \emph{sheaf of holomorphic functions} on \(\Omega\), denoted by \(\calO_\Omega\), 
        is the assignment that to each open subset \(U \subset \Omega\) assigns the ring \(\calO_\Omega(U)\) of all holomorphic functions on \(U\), 
        and set the restriction as the usual restriction of functions.
    \end{definition}

    A fundamental property of the sheaf of holomorphic functions is its coherence.

    \begin{theorem}[Oka's Coherence Theorem]\label{thm:Oka_Coherence_Theorem}
        The sheaf of holomorphic functions \(\calO_\Omega\) on an open set \(\Omega \subset \bbC^n\) is a coherent sheaf.
        \Yang{To be continued...}
    \end{theorem}

    In general, \(\calO_\Omega(U)\) is not a Noetherian ring for an open set \(U \subset \Omega\).
    However, its stalks \(\calO_{\Omega, p}\) at points \(p \in \Omega\) are Noetherian rings.
    \Yang{To be checked}

    \begin{example}\label{eg:ring_of_holomorphic_functions_is_not_Noetherian}
        \Yang{To be continued...}
    \end{example}

    \begin{proposition}\label{prop:stalk_of_sheaf_of_holomorphic_functions_is_local_noetherian_regular}
        For any point \(p \in \Omega\), the stalk \(\calO_{\Omega, p}\) of the sheaf of holomorphic functions at \(p\) is a Noetherian ring.
        \Yang{To be continued...}
    \end{proposition}

    \begin{remark}\label{rmk:topological_properties_stalk_of_sheaf_of_holomorphic_functions}
        The sheaf of holomorphic functions \(\calO_\Omega\) is a sheaf of topological rings, where the topology on \(\calO_\Omega(U)\) for an open set \(U \subset \Omega\) is given by the compact-open topology.
        \Yang{To be continued...}
    \end{remark}

    \begin{definition}\label{def:analytic_subset}
        A subset \(A \subset \Omega\) of an open set \(\Omega \subset \bbC^n\) is called an \emph{analytic subset} if for every point \(p \in \Omega\), there exists a neighborhood \(U\) of \(p\) and finitely many holomorphic functions \(f_1, \ldots, f_k \in \calO_\Omega(U)\) such that 
        \[ A \cap U = \{ z \in U : f_1(z) = f_2(z) = \cdots = f_k(z) = 0 \}. \]
        \Yang{To be continued...}
    \end{definition}
